% titel: Beispiel: Folgen
% skript: Kap. 2, S. 8

% Teil I: Werkzeuge, kein Lösungsweg.
% Regel: keine Zahlen dieser Aufgabe, keine Schrittfolge, kein Ergebnis.
% Was hier stehen soll: Begriffe, Gesetze, Gleichungen mit Gültigkeitsgrenzen,
% Vorgehen in allgemeiner Form, Einheiten, typische Fallen.

% Ein bis zwei Sätze: welche Situation, welche Frage (ohne die Lösung zu deuten).
<Situation und Fragestellung in zwei Sätzen.>

\begin{werkzeugkasten}[Werkzeuge]
\textbf{<Name des ersten Gesetzes>} in der Form, in der es im Skript steht:
\begin{equation}
  <Gleichung>
\end{equation}

\textbf{<Name des zweiten Werkzeugs>}:
\begin{equation}
  <Gleichung>
\end{equation}

\textbf{Einheiten und Bezeichner:} <welche Größe hat welche Einheit>.
\end{werkzeugkasten}

\begin{werkzeugkasten}[Vorgehen in allgemeiner Form]
<Welche Größen sind typischerweise gegeben, welche gesucht, und über welche
Beziehung hängen sie zusammen. Als Landkarte, nicht als Rezept.>
\end{werkzeugkasten}

\begin{achtung}[Typische Fallen]
\begin{itemize}
  \item <Vorzeichen, Bezugssystem, Winkel gegen Bogenmaß, gerundete Zwischenwerte ...>
  \item <Einheitenfalle, die zu dieser Aufgabe passt>
\end{itemize}
\end{achtung}
